\section{整体灵敏度与鲁棒性分析}
\label{sec:overall-sensitivity}

\subsection{核数变化与独立候选选择}
主矩阵对每个图、核数和场景独立生成候选并按官方Makespan选取最终方案，不把低核数结果强制搬运到高核数。三种场景的逐核平均速度比、五核相对四核的逐图变化见表\ref{tab:best-core-choice}。增加核数总体提高平均速度比，但有限候选搜索会使少数图的相邻核数结果出现回落。
\begin{table}[H]\centering\small
\caption{100图各核数独立选择的平均速度比及五核相对四核变化}\label{tab:best-core-choice}
\begin{tabular}{lrrrrrrr}\toprule
场景 & 1核 & 2核 & 3核 & 4核 & 5核 & 缩短/持平/变长 & 五核中位数\\\midrule
A & 1.0220 & 1.7432 & 2.3742 & 2.9520 & 3.4383 & 65/30/5 & 4.0343\\
B & 1.0564 & 1.7610 & 2.3993 & 2.9852 & 3.4738 & 70/23/7 & 4.0343\\
C & 1.0739 & 1.7998 & 2.4462 & 3.0448 & 3.5533 & 72/19/9 & 4.2035\\\bottomrule
\end{tabular}\end{table}
A、B、C的五核平均速度比分别为3.4383、3.4738和3.5533；1500条主矩阵记录均为\texttt{AI\_VERIFIED}，最终工期均不超过统一整图单核基准。五核相对四核的变长图分别为5、7和9张，原因是各核数的候选集合、候选排序和完整评价预算独立变化，表中保留这一结果以反映实际搜索行为。

\subsection{样本组成的重采样检验}
对100张图有放回抽样1500次，重新计算五核速度比均值并取2.5\%和97.5\%经验分位数。A、B、C的均值及95\%区间分别为$3.4383\,[3.1333,3.7369]$、$3.4738\,[3.1608,3.7682]$和$3.5533\,[3.2357,3.8626]$。对500条B/C配对记录重采样，$S_{\rm hw}$、$S_{\rm adapt}$和$S_{\rm opt}$的五核均值及区间分别为$1.0146\,[1.0044,1.0271]$、$1.0139\,[1.0055,1.0259]$和$1.0289\,[1.0148,1.0468]$。三段效应的区间均位于1以上，表明在当前图集合和候选范围内，Cache配置变化与C内计划调整均提供了可观测的平均工期收益。

\subsection{统一基线与受控对照}
统一基线取100张图整图单Task单核Makespan，平均为3124793.73 cycle。表\ref{tab:baseline-ablation}把三种场景五核最终选择和B/C同图同核配对分列；配对列按逐图比值取平均，严格同计划解释参见第7章。
\begin{table}[H]\centering\small
\caption{五核100图的统一基线与B/C受控对照}
\label{tab:baseline-ablation}
\setlength{\tabcolsep}{4pt}
\begin{tabularx}{\linewidth}{lYrrr}
\toprule
对照层 & 计划与配置 & 平均工期/cycle & A0加速比均值 & 配对效应比均值\\\midrule
统一基线 & 整图单Task、单核$A0$ & 3124793.73 & 1.0000 & --\\
问题一 & A场景五核最终选择 & 890650.61 & 3.4383 & --\\
问题二 & B场景五核最终选择 & 872313.38 & 3.4738 & --\\
问题三 & C场景五核最终选择 & 850429.76 & 3.5533 & --\\
B计划记录 & $F_B(X_B^*)$ & 872313.38 & -- & 1.0000\\
C初始记录 & $F_C(X_B^*)$（按配对文件） & 868575.63 & -- & 1.0146\\
C最终选择 & $F_C(X_C^*)$ & 850429.76 & -- & 1.0139\\\bottomrule
\end{tabularx}
\end{table}
500条配对记录中497条的B计划哈希与C初始计划哈希一致；case\_054的5核、case\_062的4核和case\_085的5核存在初始计划差异，差异记录保留并在配对效应中单独标记。$S_{\rm hw}$和$S_{\rm adapt}$的乘积逐条等于$S_{\rm opt}$，最大残差为$2.22\times10^{-16}$。

\subsection{依赖结构分组}
以原始计算图的弱连通分量数对100张图分层，比较五核最终选择的速度比。分组只依据输入图结构。表\ref{tab:topology-speedup}显示，分量数较多的图更容易形成并行任务，单分量图的切分收益最小。
\begin{table}[H]\centering\small
\caption{原始计算依赖分量数与五核独立选择速度比}\label{tab:topology-speedup}
\begin{tabular}{lrrrr}\toprule
原始弱连通分量 & 图数 & A均值 & B均值 & C均值\\\midrule
1个 & 16 & 1.3046 & 1.5829 & 1.5960\\
2--100个 & 57 & 3.4009 & 3.3846 & 3.4910\\
超过100个 & 27 & 4.7817 & 4.7826 & 4.8448\\\bottomrule
\end{tabular}\end{table}
单分量图的五核均值为1.30至1.60，而超过100个分量的图达到4.78至4.84。前者受长依赖链、跨核等待和结构搬运共同限制；后者可通过分组切分形成更充分的核间并行。

\subsection{候选搜索收益与实验状态}
以每个图--核数--场景的首个成功候选为起点，比较最终候选工期。A场景1至5核的平均改善为1.7900\%、38.5943\%、50.0193\%、55.5842\%和58.6697\%，对应改善图数为44、90、90、90和89；B场景为4.0928\%、39.4919\%、51.3871\%、57.2715\%和60.3622\%，改善图数为56、96、96、99和98；C场景为1.0777\%、1.3230\%、1.1666\%、0.8289\%和1.1744\%，改善图数为49、47、46、32和38。C场景的候选搜索收益较小，Cache事件时序使许多候选在同一尾部完成时刻收敛。

\begin{table}[H]\centering\small
\caption{新增实验模块及官方评价状态}\label{tab:candidate-status}
\setlength{\tabcolsep}{4pt}
\begin{tabular}{lrrrr}\toprule
模块 & 组合或图数 & 记录数 & 官方成功 & 失败记录\\\midrule
主矩阵A/B/C & 100图$\times$5核$\times$3场景 & 1500 & 1500 & 0\\
初始--最终消融 & 48个图--核--场景组 & 96 & 96 & 0\\
Cache参数固定计划 & 6图$\times$4变体 & 24 & 24 & 0\\
Cache参数重优化 & 6图$\times$4变体 & 37 & 37 & 0\\
候选排序补充 & 36个单元格 & 30 & 30 & 0\\\bottomrule
\end{tabular}\end{table}
主矩阵、消融和Cache参数实验均使用官方事件评估器；Cache rescue与容差策略属于独立诊断，不改变主矩阵的final\_selection。

\subsection{候选机制控制实验}
为检验候选排序和容量引导的作用，v7另取六张代表图，在二核和五核上比较去除单一机制后的控制程序。控制实验不改变主矩阵的候选集合，仅比较对应候选选择与主程序结果；完整评价匹配数、工期变化和计划变化均从\path{data/v7_fast_controls_20260927/controls_summary.csv}读取。
\begin{table}[H]\centering\scriptsize
\caption{v7候选机制控制实验}
\label{tab:controls-v7}
\setlength{\tabcolsep}{3pt}
\begin{tabular}{llrrrrr}\toprule
控制项 & 场景/核数 & 图数 & 匹配完整评价 & 平均相对变化/\% & 慢/平/快 & 同计划数\\\midrule
直接轻量选择 & A/2 & 6 & 2 & 0.4628 & 5/1/0 & 1\\
直接轻量选择 & A/5 & 6 & 3 & 0.8251 & 3/2/1 & 2\\
去除容量引导 & B/2 & 6 & 3 & -3.5526 & 1/4/1 & 4\\
去除容量引导 & B/5 & 6 & 0 & -- & 3/3/0 & 3\\
去除接收域 & A/2 & 6 & 2 & 0.0000 & 0/5/1 & 5\\
去除接收域 & A/5 & 6 & 3 & 0.0000 & 0/6/0 & 6\\
去除接收域 & B/2 & 6 & 3 & -0.0225 & 1/4/1 & 4\\
去除接收域 & B/5 & 6 & 0 & -- & 2/2/2 & 2\\\bottomrule
\end{tabular}
\end{table}
平均相对变化按控制工期减去主程序工期再除以主程序工期计算；负值表示控制程序在匹配样本上更快。五核去除容量引导和去除接收域的匹配完整评价数为0，表中保留这一状态，避免把未匹配的候选当作性能差异。该控制实验补充说明，接收域与容量引导共同决定候选是否进入可比较的完整评价集合。

\subsection{算法效率与评价成本}
算法成本包括图索引、候选构造与轻量排序、官方完整评价三部分。令$|V|$为非COPY操作数，$|E_0|$为原始边数，$|E|$为去重后的计算依赖数，$|\mathcal T|$为张量数，$P$为生产者--消费者配对数。索引代价为$O(|V|+|\mathcal T|+|E_0|+P)$，拓扑排序为$O((|V|+|E|)\log|V|)$，弱连通分量遍历为$O(|V|+|E|)$；就绪操作的排序在最坏情况下为$O(|V|^2\log|V|)$。若生成$M$个候选、单候选检查与评分成本为$C_{\rm cand}$，候选阶段为$O(MC_{\rm cand}+M\log M)$；完整评价随展开事件数$E_{\rm ev}$和容量扫描增长。

\begin{table}[H]\centering\small
\caption{v7主矩阵最终选择的下界距离与计时字段}
\label{tab:overall-eval-cost}
\setlength{\tabcolsep}{4pt}
\begin{tabular}{lrrrrr}\toprule
场景 & 组合数 & $F/LB$中位数 & $F/LB$第90百分位 & 评价耗时中位数/s & 评价耗时最大/s\\\midrule
A & 500 & 9.556 & 35.690 & 0.433 & 14.696\\
B & 500 & 9.252 & 35.694 & 0.666 & 21.774\\
C & 500 & 9.170 & 35.493 & 0.722 & 34.064\\\bottomrule
\end{tabular}\end{table}
表中计时字段来自主矩阵最终选择记录；未执行完整评价的基线行不计入评价耗时统计。A、B、C的中位轻量画像时间分别为0.345、0.638和0.683 s，95\%分位数分别为8.127、12.688和12.367 s。官方事件评估在CPU执行，MPS只用于轻量负载排序。

\subsection{模型评价与适用条件}
三问共用整图单核基准和接收域约束：问题一按Task接收域切分，问题二按核心接收域切分，问题三在核级方案上加入FIFO Cache查询与完成事件。问题三的三段比值恒等式、主矩阵1500条官方成功记录和Cache敏感性24/37条成功记录共同构成结果链条。候选集合有限，最终选择反映当前候选生成、排序和预算下的方案质量；算法复杂度和计时表给出了搜索开销，适用对象为依赖图、操作成本和片上容量已知且能够调用官方事件评估器的多核调度任务。
